\documentclass[11pt,letterpaper]{article}
\usepackage[utf8]{inputenc}
\usepackage[T1]{fontenc}
\usepackage[spanish,es-noshorthands,es-tabla]{babel}
\usepackage[margin=2cm,top=2.5cm,bottom=2.5cm]{geometry}
\usepackage{amsmath,amssymb}
\usepackage{enumitem}
\usepackage{fancyhdr}
\usepackage{listings}
\usepackage{xcolor}
\usepackage{tcolorbox}
\usepackage{booktabs}
\usepackage{array}
\raggedbottom
\setlength{\headheight}{14pt}
\let\vspaceoriginal\vspace
\renewcommand{\vspace}[1]{\par\vspaceoriginal{#1}\par}
\pagestyle{fancy}
\fancyhf{}
\fancyhead[L]{Basic C Without Tears}
\fancyhead[C]{\NombreDocumento}
\fancyhead[R]{Stack}
\fancyfoot[L]{Basic-C-Without-Tears}
\fancyfoot[R]{\thepage}
\renewcommand{\headrulewidth}{0.4pt}
\renewcommand{\footrulewidth}{0.4pt}
\lstdefinestyle{customc}{
  language=C,
  numbers=left,
  stepnumber=1,
  numbersep=8pt,
  numberstyle=\tiny\color{gray},
  basicstyle=\footnotesize\ttfamily,
  keywordstyle=\bfseries\color{blue!80!black},
  commentstyle=\itshape\color{green!50!black},
  stringstyle=\color{red!70!black},
  breaklines=true,
  breakatwhitespace=true,
  tabsize=4,
  showstringspaces=false,
  frame=none
}
\newtcolorbox{solucionbox}{
  before=\par,
  after=\par,
  colback=white,
  colframe=black,
  arc=0mm,
  boxrule=0.6pt,
  left=3mm,
  right=3mm,
  top=2mm,
  bottom=2mm
}
\newif\ifpauta
\ifdefined\PAUTA\pautatrue\else\pautafalse\fi

\newcommand{\NombreDocumento}{Examen 03 -- Variante B}
\begin{document}
\begin{center}
    {\LARGE \textbf{Basic C Without Tears}} \\[0.2cm]
    {\Large \textbf{Examen 03 -- Variante B}} \\[0.12cm]
    \small Fundamentos de C, matrices, control de flujo y cadenas de caracteres \\[0.08cm]
    \textbf{Autor:} Stack
\end{center}
\vspace{0.2cm}
\begin{enumerate}[label=\textbf{\arabic*.}]
\item Una matriz registra unidades producidas por tres turnos en dos máquinas:
\begin{lstlisting}[style=customc, numbers=none]
int U[3][2] = {{10, 12}, {8, 15}, {14, 9}};
\end{lstlisting}
\begin{enumerate}[label=(\alph*)]
\item Indique el valor de \texttt{U[0][1]} y \texttt{U[2][0]}.
\ifpauta\begin{solucionbox}\texttt{U[0][1]} vale 12 y \texttt{U[2][0]} vale 14.\end{solucionbox}\else\vspace{1.5cm}\fi
\item Calcule el total producido por cada turno (cada fila).
\ifpauta\begin{solucionbox}Turno 1: 22 unidades; turno 2: 23; turno 3: 23.\end{solucionbox}\else\vspace{2cm}\fi
\item Escriba código que cuente cuántos elementos son pares.
\ifpauta\begin{solucionbox}\begin{lstlisting}[style=customc]
int pares = 0;
for (int i = 0; i < 3; i++) {
    for (int j = 0; j < 2; j++) {
        if (U[i][j] % 2 == 0) {
            pares++;
        }
    }
}
printf("Pares: %d\n", pares);
\end{lstlisting}10, 12, 8 y 14 son pares; el resultado es 4.\end{solucionbox}\else\vspace{3cm}\fi
\end{enumerate}
\newpage
\item Un código de producto almacenado en \texttt{int codigo} debe estar entre 1000 y 9999, inclusive.
\begin{enumerate}[label=(\alph*)]
\item Escriba la condición de validez.
\ifpauta\begin{solucionbox}\texttt{codigo >= 1000 \&\& codigo <= 9999}.\end{solucionbox}\else\vspace{1.5cm}\fi
\item Escriba un \texttt{if-else} que imprima \texttt{Codigo valido} o \texttt{Codigo invalido}.
\ifpauta\begin{solucionbox}\begin{lstlisting}[style=customc]
if (codigo >= 1000 && codigo <= 9999) {
            printf("Codigo valido\n");
} else {
            printf("Codigo invalido\n");
}
\end{lstlisting}\end{solucionbox}\else\vspace{2cm}\fi
\item Clasifique \texttt{999}, \texttt{1000}, \texttt{4500} y \texttt{10000}.
\ifpauta\begin{solucionbox}999: inválido; 1000: válido; 4500: válido; 10000: inválido.\end{solucionbox}\else\vspace{2cm}\fi
\end{enumerate}
\newpage
\item Determine qué imprime el siguiente código.
\begin{lstlisting}[style=customc]
int M[2][3] = {{2, 1, 0}, {1, 2, 3}};
int cuenta = 0;
for (int i = 0; i < 2; i++) {
    for (int j = 0; j < 3; j++) {
        if (M[i][j] > 1) {
            cuenta++;
        }
    }
    printf("Conteo parcial: %d\\n", cuenta);
}
\end{lstlisting}
\ifpauta\begin{solucionbox}En la primera fila solo el 2 cumple; en la segunda se agregan 2 y 3. El contador acumulado vale 1 y luego 3.\par\textbf{Salida:}\\\texttt{Conteo parcial: 1}\\\texttt{Conteo parcial: 3}\end{solucionbox}\else\vspace{3cm}\fi
\item Escriba una función que cuente cuántas veces aparece el carácter \texttt{'a'} en una cadena. La función distingue mayúsculas de minúsculas.
\begin{center}\begin{tabular}{ll}\toprule\textbf{Entrada} & \textbf{Resultado}\\\midrule\texttt{"banana"} & \texttt{3}\\\texttt{"Aula"} & \texttt{1}\\\texttt{"xyz"} & \texttt{0}\\\bottomrule\end{tabular}\end{center}
\ifpauta\begin{solucionbox}\begin{lstlisting}[style=customc]
int contar_a(const char texto[]) {
    int total = 0;
    for (int i = 0; texto[i] != '\0'; i++) {
        if (texto[i] == 'a') {
            total++;
        }
    }
    return total;
}
\end{lstlisting}Se compara cada carácter exactamente con la minúscula \texttt{'a'}.\end{solucionbox}\else\vspace{3cm}\fi
\end{enumerate}
\end{document}
